\section{Related Work}
\label{sec:related-work}

% =========================================================================
% Five subsections, each one to two paragraphs of drafted prose.
% Framework is positioned as approximate Q-learning over a verifier-
% anchored context MDP, with MCTS as a policy improvement operator and
% a test-time-compute knob, not as MCTS-first.
% =========================================================================

\subsection{Search and Learning in Games}
\label{sec:rw-games}

The conceptual ancestor of the present framework is the joint use of tree
search and learned function approximation in self-play games. UCT
\citep{kocsis2006bandit} and the original Monte Carlo tree search of
\citet{coulom2006efficient} supplied the algorithmic backbone; the broader
landscape of variants and applications is surveyed by
\citet{browne2012survey}. The AlphaGo family
\citep{silver2016alphago, silver2017alphagozero, silver2018alphazero}
combined deep neural priors and value estimates with PUCT-guided search,
distilling the search policy back into the priors in a closed loop. MuZero
\citep{schrittwieser2020muzero} extended the same recipe to settings whose
dynamics must be learned. The expert-iteration formulation of
\citet{anthony2017thinking} states the underlying pattern abstractly: a
fast \textit{apprentice} policy is improved by an expensive \textit{expert}
operator (search), and the apprentice is then trained to imitate the
expert's behavior.

We inherit this joint search-and-learning structure but invert the
narrative emphasis. In the AlphaZero tradition, tree search is the central
object and the learned heads are its scaffolding; in our setting the
central object is an agent and its approximate action-value function
$\Qpsi(s, a)$, trained against verifier-anchored Bellman targets, and
search enters as one of several mechanisms that supply stronger training
targets and an optional test-time-compute knob
(Sections~\ref{sec:method-training} and~\ref{sec:method-inference}). The
shift is what is appropriate when (a) the environment is stochastic and
LLM-induced rather than a deterministic two-player game, so a chance-node
estimator becomes essential and a collapsed single-node tree becomes
biased (Section~\ref{sec:method-mcts}); and (b) the value channel is
anchored by a non-LLM verifier rather than by game outcomes alone, which
makes a wider set of non-search exploration policies viable as baselines
and makes the choice between $\epsilon$-greedy, softmax, and search-driven
data collection an empirical rather than an architectural question.

\subsection{Verifier-Anchored LLM Reasoning}
\label{sec:rw-verifiers}

Several recent systems share the principle that learning signal for
reasoning models should come from non-LLM verifiers rather than from
human-rated aesthetic critics. \citet{lightman2024verify} demonstrated
that process reward models trained on per-step labels for MATH solutions
outperform outcome-only supervision on the same benchmark, and
\citet{cobbe2021gsm8k} introduced GSM8K together with the original
verifier-style training pipeline for grade-school arithmetic. AlphaCode
\citep{li2022alphacode} used unit-test execution as the verifier for
competitive programming, and AlphaProof \citep{deepmind2024alphaproof}
extended the same architectural pattern to formal mathematics with Lean
as the verifier. \citet{guo2025deepseekr1} report large-scale
reinforcement learning against rule-based verifiers for math and code
without a learned reward model, deliberately avoiding the Goodhart
attack surface that attends preference-trained critics.

Our work shares the verifier-anchored reward principle but contributes
along two axes. The first is framework-level: Section~\ref{sec:method-reward}
makes the verifier-versus-aesthetic asymmetry a load-bearing design
commitment rather than an implementation detail. Verifier outcomes are the
only signals allowed to enter the Bellman target; learned aesthetic
signals such as reasoning-quality classifiers or deception detectors may
enter only as MCTS priors, where they bias exploration but cannot rewrite
the value function. The second is architecture-level:
$\textsc{InvokeVerifier}$ is a callable verb in the agent's action space
(Section~\ref{sec:method-actions}) rather than a post-hoc parser pass
applied only at the terminal answer. The agent decides when to spend a
verifier call, intermediate verifier outputs land in the structured state
and become available to the heads, and a hypothesis can be promoted to a
verified fact mid-trajectory. The choice converts verifier interaction
from a wrapper around the model into a first-class part of the
deliberation, which is what makes verifier-anchored intermediate rewards
(Equation~\ref{eq:rewardterminal} and the fallback discussion that
follows it) consistent with the no-Goodhart principle.

\subsection{Test-Time Compute Scaling}
\label{sec:rw-ttcs}

\citet{snell2024scaling} established empirical scaling laws relating
test-time compute to accuracy on verifiable benchmarks, showing that
optimal inference-time allocation can match or exceed the returns of
proportional pretraining-scale increases. OpenAI's o1 line
\citep{openai2024o1} is the public anchor for test-time-compute-as-product
in deployed reasoning systems. A common pattern in this literature is the
inference-time technique (chain-of-thought \citep{wei2022cot},
self-consistency, best-of-$N$, beam search over rationales) studied as a
post-hoc augmentation of an already-trained generator.

Our contribution is methodological rather than algorithmic in this
register. Section~\ref{sec:exp-baselines} pre-registers a baseline grid
that separates fast-mode policy improvement ($B = 0$, where the trained
$Q$-head is acted on greedily) from search-mode improvement ($B > 0$,
where MCTS refines the head estimate at the query), with matched
inference compute across all conditions. The separation matters because
the literature has often conflated two distinct claims: ``the model is
better at compute $X$'' (which can be true purely because the trained
heads are stronger) and ``the model uses additional compute $X$
usefully'' (which is the genuine test-time-compute question). By writing
the compute axis, the baseline grid, and the headline frontier figure
before any training begins, we make the conflation difficult to commit
in either direction. The same evaluation protocol turns the search budget
$B$ from a tuned hyperparameter into a continuous dial the user
controls at inference, which is the practically relevant form of a
test-time-compute knob.

\subsection{MCTS for Language Model Reasoning}
\label{sec:rw-mcts-llm}

A growing body of recent work applies tree search to language model
reasoning. \citet{feng2023alphazerollm} adapt the AlphaZero recipe
directly, using a learned value model and PUCT to guide token-level or
step-level decoding. \citet{xie2024monte} use MCTS to generate preference
data for iterative DPO-style training, treating search as a trajectory
generator whose output is consumed downstream. The closest prior
published system to our framework is rStar-Math
\citep{guan2025rstarmath}, which combines small base language models
with MCTS-based self-training and a process reward model trained from
MCTS rollouts; the broad approach (small LMs, search-guided rollouts,
distillation from search statistics) is the same family our framework
inhabits, and rStar-Math's empirical results establish that the family
works on competition-grade math. The general shape of these systems is
recognizable from the AlphaZero ancestor: tree search produces better
trajectories than the base policy, and the base policy is updated toward
what search would have done.

Our work differs along four structural axes that together distinguish a
framework paper from a fitted system. First, the search tree explicitly
bifurcates decision and chance nodes (Section~\ref{sec:method-mcts}); this
is what makes the chance-node mean $W(s, a) / N(s, a)$ a legitimate
estimator of the inner Bellman expectation when the environment is
stochastic, as it is whenever the transition is induced by LLM sampling
rather than by deterministic game rules. Collapsing the two node types,
the pattern most often seen in deterministic-game tutorials, biases every
ancestor's $Q$-estimate. Second, the agent acts over a small fixed set of
coarse verbs $\Aset$ (Section~\ref{sec:method-actions}) and the
$Q$-function is a multi-head categorical output rather than a scorer over
$(s, a)$ embeddings of free-form text. This trades expressivity for a
stable categorical distribution that PUCT can branch on and that the
policy head can match by cross-entropy. Third, $\textsc{InvokeVerifier}$
is one of the verbs the agent may take rather than a step performed only
on the committed answer (Section~\ref{sec:method-actions}); the verifier
is a learnable resource the policy budgets. Fourth, the architecture,
action space, evaluation protocol, and falsifiable hypotheses are
pre-registered (Section~\ref{sec:experimental-plan}) before any training
begins, which is what we take to distinguish a methodology paper from a
post-hoc justification of an empirically tuned system.

Two of these axes deserve specific contrast with
\citet{guan2025rstarmath}. On the state and action axis, rStar-Math
operates over token-level or step-level decoding, with search nodes
keyed to partial generations of natural-language reasoning; our context
MDP is a structured-state object with a fixed seven-verb action space
and id-based targeting (Sections~\ref{sec:method-state}
and~\ref{sec:method-actions}), so the categorical action set is small
and stable across problems and chance nodes can be addressed by
identifier rather than by string match. On the methodology axis,
rStar-Math is a fitted system whose architectural choices, reward
shaping, and evaluation protocol are reported alongside its empirical
results; our framework pre-registers the baseline grid, statistical
analysis plan, and falsifiable hypotheses (Section~\ref{sec:experimental-plan})
before training begins. The two efforts are complementary: rStar-Math
demonstrates empirically that small base LMs with MCTS-driven
self-training reach competitive math accuracy, and our framework aims
to make the design space and the falsification criteria of that broad
approach explicit.

\subsection{Self-Improvement and Synthetic Trajectory Generation}
\label{sec:rw-selfimprove}

A complementary line of work asks whether a language model can improve
itself by generating its own training data. STaR \citep{zelikman2022star}
samples rationales from the base model, retains those whose final answer
matches a known label, and fine-tunes on the retained set; the procedure
is iterated. V-STaR \citep{hosseini2024vstar} trains a verifier on the
rationales STaR discards and uses it at inference to rerank candidates.
Both approaches operate at the granularity of complete trajectories: the
unit of credit is the whole solution path, and the supervision signal is
whether the path's final answer is correct.

Our framework sits in the same family of search-then-distill systems but
distills at a finer granularity. MCTS produces locally improved
trajectories, and the heads are trained on per-step chance-node estimates
$\Qmcts(s, a)$ rather than only on terminal correctness
(Section~\ref{sec:method-training}). What makes per-step distillation
feasible is the structured intermediate state of
Section~\ref{sec:method-state}: a math state with explicit
\texttt{Fact}, \texttt{Subgoal}, \texttt{Hypothesis}, and verifier-call
records is a Markov state with a well-defined value, whereas a raw
token-prefix has neither. The per-step credit assignment is closer to
classical $Q$-learning than to trajectory-level rejection sampling: the
chance-node $Q$-estimate at $(s, a)$ averages over the discounted return
of every continuation that branches from that action, so a correct
intermediate hypothesis that leads to an incorrect final commit still
receives a positive partial-credit signal at the chance node where it was
introduced. The structural-state choice is therefore what allows the
search-then-distill pattern of STaR's lineage to be carried over into a
Bellman-style update rather than a trajectory-level supervised target.
